\documentclass[aps, pre, reprint]{revtex4-2}
\usepackage[utf8]{inputenc}
\usepackage{orcidlink}
\usepackage{natbib}
\usepackage{amsmath}
\usepackage{hyperref}
\usepackage{graphicx}

% Fig 1:

% Schematic of networkt
% Schematic of 2D 
% Picture of lattice in 2D
% Histogram of colours and numbers in 2D (?)

% Fig 2:

% Mix/max spot size over time
% Bottom: histogram of emissions

% Fig 3:

% Spot size and emission histograms over time. 
% In grey, plot random mixed, and also in different colours plot different density results

\begin{document}

\title{A spectrum of power laws in a SOC model for coronal mass ejections}
\author{Riz Fernando Noronha \orcidlink{0009-0007-2923-3835}, Kim Sneppen \orcidlink{0000-0001-9820-3567}}
\email{noronha@nbi.ku.dk}
\affiliation{Niels Bohr Institute, University of Copenhagen, Copenhagen 2200, Denmark}

\begin{abstract}
    Solar flares and coronary mass ejections are phenomena that emit energy or mass from a star. This process has been modelled in the past through the merging and creation of magnetic loops to create power-law distributions of emission sizes, however, the emission exponent is far too step ($\tau\approx$3, in contrast with experimental evidence of $\tau\approx1.5-1.6$). We modify the above model by changing the dynamics by which new magnetic loops are injected into the system, such that their poles typically neighbor poles of the same sign. The critical exponent of the modified model is reduced to $\tau\approx 1.9$, better matching real-world data.
\end{abstract}


\maketitle



\section{Introduction}



The area beneath the surface of the sun is home to a strong magnetic field, which varies across the surface. Regions with particularly high magnetic fields typically expel magnetic flux tubes into the solar atmosphere. As per a model described by Parker \cite{ParkerSolarFlares}, if local magnetic field gradients become sufficiently steep, the coronal magnetic field can change it's topology through reconnecting some of the loops, thus releasing energy. \\

Data gathered from energy releases during solar flares appears to follow a power-law distribution \cite{NishizukaXRay, AschwandenSolarFlareData}. This seems to suggest that the surface of the sun is in a state of self-organized criticality. In order to investigate this system further, several models have been established (such as Lu and Hamilton's model \cite{LuHamiltonSandpile}, which is essentially a sandpile model similar to the one derived by Bak, Tand and Wiesenfeld \cite{BakSandpile}), but for the purposes of this manuscript we discuss the model proposed by Hughes et al. \cite{HughesSolarFlares} and Sneppen et al \cite{sneppenTrusinaMeanFieldBipartite}. \\


\section{The Model}


The model works with the general concept of combining and reconnecting nodes in a network. We can describe the regions of high magnetic field on the suns surface (typically known as `active regions') by nodes, some of them being positive nodes and some being negative nodes. We can then make links connecting every positive node to a negative one. Now, every time step, we perform the following two steps:
\begin{itemize}
    \item Merge two nodes: We choose two random nodes, and combine them. The effects change depending on whether both nodes have the same sign or not: if they do, say one has 2 positive links and the other 3 positive links, we simply combine them to create a positive node with 5 links (with links to the same nodes it's components were linked to). However, if the two nodes have different signs, the merged node has the same sign as the larger of the two, and all the links of the smaller node are directly reconnected to available links of the larger node (\autoref{fig:bipartite}).
    \item Inject new nodes: As the combining of nodes will gradually decrease the number of nodes in the system, nodes connected to each other (i.e, loops) are injected at a fixed rate.
\end{itemize}

\begin{figure}[h!]
    \centering
    \includegraphics[width=\linewidth]{fig/bipartite.png}
    \caption{Merging of nodes in a bipartite network, and how connections are dealt with. Left: Merging nodes of the same sign, Right: Merging nodes of different signs \cite{ComplexLectureNotes}.}
    \label{fig:bipartite}
\end{figure}

Note that this entire bipartite model can be simulated without having to store data about which nodes are linked to which, as it is possible to only store the number of links in each node and adjust for merging accordingly. The model discussed in the paper has several more complications: it is run on a 2D grid, and allows nodes to reconnect (switch) loops if within a certain distance of one another, with reflective boundary conditions (such that the nodes are elastically reflected back at the boundaries). It has several free parameters, such as $m$, the rate at which new loops are injected in (which determines the cutoff), $l_{nl}$, the distance between the poles of a newly injected loop, $l_{min}$, a length scale below which loops agglomerate, etc. \\

When two nodes of the opposite sign combine, the resulting loss of links is treated as an \textit{emission}. That is, if two nodes of $+4$ and $-1$ combine, the resulting node has a sign of $+3$ and an emission of $1$ is recorded. The emission distributions followed a power law with a critical exponent of $3$, while the distribution of the connectivity (regardless of the sign) of nodes had an exponent of $2$.


\section{The Model}

The power laws obtained by the previous model were too steep: we observe larger amounts of extreme events than the model predicts. Nishizuka et al. \cite{NishizukaXRay} report a power law of $\tau=1.5$ in emission intensity, while deArcangelis \cite{DeArcangelisSolarFlareData} report an exponent of $\tau=1.65$, both far lower exponents than the $\tau=3$ that Hughes' model predicts. Wang and Dai \cite{WangXRayData} also find exponents of $1.53$ and $1.65$ depending on the type of solar flare. The goal of this project is to take a simplified form of the above model (with less free parameters) and modify it to try to decrease the critical exponent, thus getting results more similar to what is observed in nature. This is done by changing the rules of injecting new loops in the system: rather than injecting them randomly, we attempt to preferentially inject them next to neighbours of a similar polarity.


\subsection{The 2-Dimensional Model}

A two dimensional $L\times L$ matrix is generated as a representation of the solar surface. A density parameter $\rho$ determines the spareseness of the matrix. In particular, the number of occupied cells in the matrix, $N$, is calculated by
\begin{equation}
    N = 2\left\lfloor\frac{\rho L^2}{2}\right\rfloor
\end{equation}
By the above definition, $N$ will always be an even integer. $N$ random cells are selected and populated with non-zero values, initialized to be $+1$ and $-1$. Half of the cells are set to $+1$, and the other half to $-1$, so as to ensure that the sum of all entries in the matrix is 0. \\

Every non-zero point on the matrix then undergoes a random walk, with periodic boundary conditions. This is done in a monte-carlo simulation by:
\begin{itemize}
    \item Choosing a particular `occupied' matrix cell
    \item Choosing one of the four nearest neighbours to move to
\end{itemize}
If the chosen destination is already occupied with a non-zero number, it's value is updated to be the sum of the original and the newly moved number. In case the two numbers have opposing signs, the absolute difference between them is treated as an `emission'. \\

In order to conserve the total number of occupied cells, every time the number of active points drops below $\left(N-1\right)$, we introduce two more: a positive ($+1$) and a negative $(-1)$, in a random vacant cell. Alternatively, in our local model we generate the new poles close to existing ones: the new positive node $(+1)$ is introduced in a cell which neighbours am existing positive node, and vice versa for the new negative node. \\

It is important to note that (especially at high densities), a possible location for the new nodes might not exist. For example, at $\rho=1$ you will only have 2 vacant cells, and if both of them are surrounded by only negative nodes, then you cannot place the positive node. For the sake of the project, this phenomenon was not investigated (as it involves changing the injection rules), and the only systems shown are those which never ran into this problem. A short note on this can be seen in the \autoref{sec-discussion}.

\section{Methods}

The simulations were run on the following models for 200 million time steps (where a single node is moved, and nodes injected if necessary every time step), where the first 100 million timesteps were assumed to contain the approach to equilibrium, and data was only collected in the latter 100 million time steps. The following data was collected:

\begin{itemize}
    \item The state: the configuration of the lattice was stored every $100$ iterations. This data was not recorded every step, as there would be a substantial overlap between adjacent time steps (as only one spot can move in each time-step). Recording data every time step will not be necessarily wrong, but it will greatly increase the amount of data you have to work with for a diminishing reward. The number $100$ was chosen fairly arbitrarily.
    \item The spot sizes: Every $100$ iterations, we use the state to calculate the number of nodes on the lattice with a particular connectivity, regardless of their sign (positive or negative). If one were to link the number of loops attached to a particular node with a notion of a `size', this can be thought of as recording a histogram of the sunspot sizes. As the spots can have very large sizes (after all, they stem from a power-law distribution, which can have extreme events), the data is log-binned, and then divided by the number of discrete integers in each bin.
    \item The emissions: Every step, if two nodes happen to merge, the emission is calculated as
    $$\min\left(\,\lvert p_1\rvert, \lvert p_2\rvert\,\right)$$
    where $p_1$ and $p_2$ are the values of the two nodes. The emission data is recorded, and log-binned in a similar manner to the spot-size distribution described above.
\end{itemize}

Using the state data, the number of domain walls for the 1-dimensional simulations is also calculated.


\section{Results}


\autoref{fig:2demissions} below contains plots for the two statistics in the two-dimensional case. It can be seen that the power laws once again correspondingly decrease: in the case of the sunspot sizes, the change to neighbour injections decrease the critical exponent from $2$ to $1.5$, and for the emissions, from $3$ to $2$. While the emission exponent doesn't match the data as well as the one-dimensional case, it is still a significant improvement from the extremely steep exponent of $3$ in the prior (random injection) model.

\begin{figure}[h!]
    \centering
    \includegraphics[width=0.49\linewidth]{fig/2d_spotsize.png}
    \includegraphics[width=0.49\linewidth]{fig/2d_emissions.png}
    \caption{Left: Histograms of the spot sizes, Right: Histograms of the emission sizes, for the 2-dimensional model. Both the random injection and neighbour injection models are shown. Both models were set to have $L=64, \rho=0.2$.}
    \label{fig:2demissions}
\end{figure}

Both the density $\rho$ (Appendix \autoref{fig:2drho}) and the system size $L$ (Appendix \autoref{fig:2dL}) had a clear impact on the cutoff. At first glance, it appears that they have a slight impact on the critical exponent of the emission data as well, however, it could be that we simply see the cutoff behaviour dominate and drag the curve down at those points. The later points have extremely low statistics, and critical systems such as these tend to have a certain level of robustness in their results. Furthermore, if the distribution of node sizes is unchanged, the distribution of emissions should not change, as the two are clearly linked (as the formation of a large emission necessarily requires two big nodes of opposite polarity to exist). \\


\section{Discussion}\label{sec-discussion}

With higher densities or longer simulations, it is likely that one runs into the problem of not having an appropriate location to initialize new spots. In this case, the most intuitive answer would simply be to first test if they are eligible locations, and if there aren't, choose a random location to initialize a spot in. However, though this seems like a very simple and natural extension of the current simulation, in practise it complicates the matter. What we're doing here is changing the injection rules, which will have an impact on the system by changing the time-evolution of the order parameter. \\

In the 1D case, this new method of injections makes it possible for domain walls to be created, as opposed to the previous case where they could only be destroyed and the system necessarily coarsened with time. Furthermore, this `new rule' of choosing a random location will not be triggered at a constant rate, but rather one that depends on density: for high densities, we would expect it to happen fairly often while it's an extremely rare occurrence for low densities. This makes the density a parameter of interest, harming the universality of the solution. \\

However, there is also an advantage: in the extreme case where we get unlucky every time and a random neighbour is picked, we would simply have the random initialization model, which would give us a critical exponent of $3$ in the emissions. In the other extreme, we simply have the neighbour initialization model described in the paper, with a critical emission exponent of $1.7$ (in 1 dimension). Thus, by tweaking the density parameter, we can change the system to give us power laws with exponents lying somewhere between $1.7$ and $3$ (though possibly not exhaustive, as even at maximum density we would not get a power law of $2$: however, it would definitely be greater than $1.7$ because of the new effect). This scaling of the critical exponents with $\rho$ can be considered a Griffith's Phase, first discussed by Griffiths in the context of a diluted Ising model \cite{griffithsPhase}. 

The waiting-time between subsequent solar flares appears to follow a time-dependent Poisson distribution \cite{WheatlandWTD}. However, over long periods, the distribution appears to behave as a power-law \cite{BoffettaSolarFlares}. In isolation, the sandpile SOC model was unable to replicate the power-law behaviour \cite{MayaBTWWaitingTimes}, as the avalanches are temporally uncorrelated. Investigating the waiting time distribution of the model described above could be yet another avenue for exploration. \\

In the 2D model, changing the density and system size appeared to have a marginal effect on the critical exponent (Appendix \autoref{fig:2drho} and \autoref{fig:2dL}). Further simulations run on more data could clarify as to whether this is an actual effect or a statistical problem of too small simulations. One potential justification for it being an actual effect is to consider a site chosen as an injection location for a new positive node: the only restriction is that the site has to neighbour at least one positive node. At low densities, the other neighbours will probably be empty, but at high densities they're more likely to be occupied. If the occupied neighbours also include negative nodes, then the new node is no closer to a positive node than a negative one, and so in the next step it might either agglomorate positively or negatively. This can also explain why we don't see this effect in 1D, as there we have only 1 potential neighbour which could have a negative node, while we have 3 other nearest neighbours in 2D. \\

In addition to its relevance relating to solar flares, the model we describe could potentially be used in a completely mathematical context as a novel algorithm to generate scale-free bipartite networks with the given exponents.


\section{Conclusion}

In this paper, we demonstrate that by making a simple modification an already establish SOC model for solar flare dynamics, we can significantly modulate the critical exponents of the system. By preferentially injecting new nodes next to existing ones of the same polarity, the critical exponents can be dropped fairly effectively (the emission exponent drops from $3$ to $1.9$). This allows the model to correspond much better with real-world data (which had an exponent of $\tau\approx1.5-1.65$). \\

Further investigations into the waiting-time distribution between emissions or modifying the model further to obtain an arbitrary critical exponent could provide more insight into the system's underlying dynamics.


\bibliography{biblio}


\clearpage

\appendix

\section{Mean-field renormalization}

We neglect the role of spatial structure, and focus on a mean-field approximation. Here, any two spots randomly combine each timestep. This can lead to two cases:
\begin{enumerate}
    \item Same signs: Two spots of the same sign merge, i.e, $i,j\to i+j$.
    \item Different signs: Two spots annihilate, leaving a spot of size $\lvert i-j\rvert$, and an emission of $\min(i,j)$.
\end{enumerate}
We define the quantity $q$ as the probability that a collision has the same sign. 
This correlates with the parameter $p$ from the lattice model: having purely random injections leads to merging and annihilation being equally likely ($q\!=\!0.5$), however preferentially injecting near one's neighbours leads to merging being more likely ($q\!\to\!1$). $q$ would thus be a monotonic function of $p$.

Let $n_m$ be the number density of positive spots of size $m$, and $N=\sum_{m=1}^\infty m$ be the total number density. By symmetry, the negative spots will have the same distribution.

A spot of size $m$ changes size as soon as it collides with anything (regardless of the sign). Thus, the loss of sites of size $m$ is given by

\begin{equation}
    \textrm{loss} = n_m N
\end{equation}

The number of spots of size $m$ can increase in two ways. The first is by two smaller spots of the same sign combining into one. As collisions of the same sign are biased by $q$, we have

\begin{equation}
    \textrm{gain}^+ = \frac{q}2 \sum_{i+j=m} n_i n_j
\end{equation}

The factor $1/2$ comes due to double counting in the sum, as $(i,\!j)=(j,\!i)$.

In addition, a spot larger than $m$ could collide with a spot of the opposite sign, and thus shrink to size $m$:

\begin{equation}
    \textrm{gain}^- = (1-q) \sum_{j\geq1} n_{m+j} n_j
\end{equation}

In the steady state, gain=loss, and thus:

\begin{equation} \label{eq:meanfield-gain-loss}
    \frac{q}{2}\sum_{i+j=m} n_i n_j + (1-q)\sum_{j\ge1} n_{m+j} n_j = n_m N
\end{equation}

It should be noted that this equation is only true for $m\geq2$, as $m=1$ contains injections to account for the loss of mass through emissions. \autoref{eq:meanfield-gain-loss} has no explicit scale dependence, and is invariant to rescaling the mass from $m\to bm$. As such, the solution lends itself to a power law form, $n_m \sim m^{-\tau}$.

\section{Renormalization}

To analyze the exponent, we consider the continuum case. We define two collision integrals to be

\begin{align}
    C(\mu) &= \int_1^{\mu-1} n(i)\,n(\mu-i)\,di \\
    K(\mu) &= \int_1^{\infty} n(\mu+j)\,n(j)\,dj
\end{align}

Where $n(i)\sim i^{-\tau}$, in which we assume $1<\tau\leq 2$. \autoref{eq:meanfield-gain-loss} now reads

\begin{equation} \label{eq:renormalization-gain-loss}
    \frac{q}{2}C(\mu) + (1-q)K(\mu) = n(\mu) N
\end{equation}

In addition, $N$ can be calculated to be

\begin{align}
    N &= \int_1^\infty A m^{-\tau} dm \\
    &= \frac{A}{\tau-1}
\end{align}

Which is finite for $\tau > 1$.

We then rescale $i$ by a factor $\mu$, so $i \!=\! \mu x$ and $j \!=\! \mu y$, and define $\delta=1/\mu$. This gives us

\begin{align}
    C &= A^2\mu^{1-2\tau}\!\!\int_\delta^{1-\delta}\!\! x^{-\tau}(1-x)^{-\tau}dx \\
    K &= A^2\mu^{1-2\tau}\!\!\int_\delta^{\infty}\!\! y^{-\tau}(1+y)^{-\tau}dy .
\end{align}

Both integrals diverge as $\delta\to0$. In $C$, as $x\to0$ only the factor $x^{-\tau}$ misbehaves, while $(1-x)^{-\tau}\to1$. Expanding the well-behaved factor and multiplying out,

\begin{equation}
    x^{-\tau}(1-x)^{-\tau} = x^{-\tau} + \tau\,x^{1-\tau} + \tfrac{\tau(\tau+1)}{2}\,x^{2-\tau} + \dots
\end{equation}

Integrating term by term from $\delta$, a term $x^{-s}$ survives the limit only if $s<1$. The first term has $s=\tau>1$ and diverges, while the second has $s=\tau-1<1$ and converges, for $\tau\leq 2$.

$C$ diverges at both ends, $x=0$ and $x=1$. $K$ diverges at $y=0$ by the same argument, but not at $y\to\infty$. 

We isolate the divergences by adding and subtracting the misbehaving powers:

\begin{align}
    x^{-\tau}(1\!-\!x)^{-\tau} &= \underbrace{\Big[x^{-\tau}(1\!-\!x)^{-\tau}-x^{-\tau}-(1\!-\!x)^{-\tau}\Big]}_{\sim\tau x^{1-\tau}-1\ \text{as } x\to0} + x^{-\tau} + (1\!-\!x)^{-\tau} \\
    y^{-\tau}(1\!+\!y)^{-\tau} &= \underbrace{y^{-\tau}\Big[(1\!+\!y)^{-\tau}-1\Big]}_{\sim-\tau y^{1-\tau}\ \text{as } y\to0,\ \ \sim -y^{-\tau}\ \text{as } y\to\infty} + y^{-\tau}
\end{align}


Each bracket is integrable over the full range for $1<\tau\leq2$, and the leftover pure powers are elementary, $\int_\delta^{1-\delta}x^{-\tau}dx=[\delta^{1-\tau}-(1-\delta)^{1-\tau}]/(\tau-1)$ and $\int_\delta^\infty y^{-\tau}dy=\delta^{1-\tau}/(\tau-1)$. Hence

\begin{align}
    \int_\delta^{1-\delta}\!\!\! x^{-\tau}(1-x)^{-\tau}dx &= \frac{2\delta^{1-\tau}}{\tau-1}+B_c+O(\delta^{2-\tau}) \\
    \int_\delta^{\infty}\!\!\! y^{-\tau}(1+y)^{-\tau}dy &= \frac{\delta^{1-\tau}}{\tau-1}+B_a+O(\delta^{2-\tau}) 
\end{align}

with the subtracted constants

\begin{align}
    B_c&=\int_0^1\!\Big[x^{-\tau}(1-x)^{-\tau}-x^{-\tau}-(1-x)^{-\tau}\Big]dx-\frac{2}{\tau-1} \\
    B_a&=\int_0^\infty\!\! y^{-\tau}\Big[(1+y)^{-\tau}-1\Big]dy
\end{align}

Both constants reduce to Beta functions, by different routes.

For $B_c$, when $\tau<1$ the two subtracted powers are separately integrable. $\int_0^1x^{-\tau}dx=\int_0^1(1-x)^{-\tau}dx=1/(1-\tau)$, and thus together they cancel the $-2/(\tau-1)$ term. $B_c$ is thus the unsubtracted integral, which is the first Euler form

\begin{align}
    \int_0^1 x^{a-1}(1-x)^{b-1}dx &= B(a,b) \\
    &= \frac{\Gamma(a)\Gamma(b)}{\Gamma(a+b)}
\end{align}

with $a-1=b-1=-\tau$, which gives us $B_c=\Gamma(1-\tau)^2/\Gamma(2-2\tau)$. The subtracted definition converges for all $\tau<2$ and this expression is analytic there. Thus, the agreement on $0<\tau<1$, is valid on $1<\tau<2$.

The previous route is unavailable for $B_c$, as $y\to\infty$ the bracket tends to $-1$ and the integrand to $-y^{-\tau}$, so $B_a$ exists only for $\tau>1$. Instead, we integrate by parts, with $u=(1+y)^{-\tau}-1$, $dv=y^{-\tau}dy$. The boundary term $[uv]_0^\infty$ vanishes at both ends: as $y\to0$,
$u\simeq-\tau y$ and $uv\sim y^{2-\tau}\to0$ since $\tau<2$; as $y\to\infty$,
$u\to-1$ and $v\sim y^{1-\tau}\to0$ since $\tau>1$. Hence

\begin{align}
    B_a &= -\int_0^\infty v\,du \\
    &= \frac{\tau}{1-\tau}\int_0^\infty y^{1-\tau}(1+y)^{-\tau-1}dy
\end{align}

The remaining integral is the second Euler form $\int_0^\infty y^{a-1}(1+y)^{-a-b}dy=B(a,b)$, with $a-1=1-\tau$ and $a+b=\tau+1$, i.e. $a=2-\tau$ and $b=2\tau-1$. As both are positive on $1<\tau<2$, it converges outright, with no continuation necessary. Using $\Gamma(2-\tau)=(1-\tau)\Gamma(1-\tau)$ and $\Gamma(\tau+1)=\tau\Gamma(\tau)$, the prefactor cancels:

\begin{align}
    B_a &= \frac{\tau}{1-\tau}\cdot\frac{\Gamma(2-\tau)\Gamma(2\tau-1)}{\Gamma(\tau+1)} \\
    &= \frac{\tau}{1-\tau}\cdot\frac{(1-\tau)\Gamma(1-\tau)\Gamma(2\tau-1)}{\tau\,\Gamma(\tau)} \\
    &= \frac{\Gamma(1-\tau)\Gamma(2\tau-1)}{\Gamma(\tau)} .
\end{align}


Putting them together, we have

\begin{align}
    B_c(\tau)&=\frac{\Gamma(1-\tau)^2}{\Gamma(2-2\tau)} \\
    B_a(\tau)&=\frac{\Gamma(1-\tau)\,\Gamma(2\tau-1)}{\Gamma(\tau)}
\end{align}

Using $\delta^{1-\tau}\mu^{1-2\tau}=\mu^{-\tau}$ and $N=A/(\tau-1)$,

\begin{align}
    C(\mu)&=2n(\mu)N + A^2B_c\,\mu^{1-2\tau} \\
    K(\mu)&=n(\mu)N + A^2B_a\,\mu^{1-2\tau}
\end{align}

Substituting these into  \autoref{eq:renormalization-gain-loss}, the $n(\mu)N$ terms cancel the right-hand side and we get

\begin{equation}
    \frac{q}{2}B_c(\tau) + (1-q)B_a(\tau) = 0 
\end{equation}

Which can then be reduced to 

\begin{align}
    \frac{q}{1-q} &= - \frac{2B_a}{B_c} \\
    &= -2\,\frac{\Gamma(2\tau-1)\,\Gamma(2-2\tau)}{\Gamma(\tau)\,\Gamma(1-\tau)}
\end{align}

By using Euler's reflection formula, $\Gamma(z)\,\Gamma(1-z)=\pi/\sin\pi z$, we can derive

\begin{equation}
    \cos(\pi \tau) = \frac{1-q}q
\end{equation}

Although this has infinitely many solutions, for $\tau \in (1, 2]$, $\cos(\pi\tau)$ increases monotonically from $-1$ to $1$, and as such there is only one root. We can solve for $\tau$, and finally obtain

\begin{equation}
    \tau(q) = 2 - \frac1\pi \arccos\left(\frac{1-q}{q}\right)
\end{equation}





\end{document}

